% Compiler: pdfLaTeX
\documentclass[11pt,a4paper]{article}

\usepackage[utf8]{inputenc}
\usepackage[T1]{fontenc}
\usepackage{lmodern}
\usepackage{microtype}
\usepackage{geometry}
\usepackage{enumitem}
\usepackage{titlesec}
\usepackage{fancyhdr}
\usepackage{hyperref}

% Document details
\newcommand{\researchtitle}{The planning fallacy}
\newcommand{\researchshorttitle}{The planning fallacy}
\newcommand{\researchdate}{2026-10-27}

% Layout
\geometry{left=2.7cm,right=2.7cm,top=2.5cm,bottom=2.5cm,headheight=14pt}
\hypersetup{hidelinks,pdftitle={\researchtitle}}
\urlstyle{same}
\frenchspacing
\setlength{\parindent}{0pt}
\setlength{\parskip}{0.55em}
\linespread{1.04}
\setlength{\emergencystretch}{1em}

\titleformat{\section}
    {\large\bfseries}{\thesection}{0.65em}{}
\titleformat{\subsection}
    {\normalsize\bfseries}{\thesubsection}{0.65em}{}
\titlespacing*{\section}{0pt}{1.15em}{0.35em}
\titlespacing*{\subsection}{0pt}{0.7em}{0.2em}
\setlist[itemize]{leftmargin=1.35em,itemsep=0.25em,topsep=0.2em}
\setlist[enumerate]{leftmargin=1.6em,itemsep=0.25em,topsep=0.2em}

\pagestyle{fancy}
\fancyhf{}
\fancyhead[L]{\small Random Research}
\fancyhead[R]{\small\researchshorttitle}
\fancyfoot[C]{\small\thepage}
\renewcommand{\headrulewidth}{0pt}
\renewcommand{\footrulewidth}{0pt}

\begin{document}

{\LARGE\bfseries\researchtitle\par}
\vspace{0.4em}
{\small\researchdate\par}

\section{Overview}

The planning fallacy is the tendency to underestimate how long a task
will take, even when similar tasks have taken longer before.
Buehler, Griffin, and Ross studied this pattern in everyday academic
and nonacademic work. \cite{buehler}

A plan can sound reasonable because its steps are easy to imagine.
The difficulty is deciding how closely those steps resemble the way
work actually gets finished.

\section{Findings}

In Buehler and colleagues' studies, participants gave overly optimistic
estimates of their own completion times. When explaining their forecasts,
they focused mainly on how the upcoming task would unfold. They also
tended to explain previous delays through temporary or specific
circumstances, making those experiences seem less relevant to the
next attempt. \cite{buehler}

In one study, asking participants to connect relevant past experiences
with their predictions eliminated the optimistic bias. That is a useful
result within the study, though it does not establish a correction
that will work for every task. \cite{buehler}

Kahneman and Lovallo describe an inside view that builds forecasts around
the current plan and its scenarios. Their account explains how attention
to the particular project can crowd out information about previous
outcomes. \cite{kahneman}

\section{Example}

Suppose someone sets aside an evening to write a short research note.
They allow time for reading and drafting, then expect to be finished.

While reading, they find that two sources use the central term differently.
They spend time resolving the difference, replace an example, and revise
the draft. None of these steps is especially surprising once it happens,
but the initial estimate left little room for them.

This is a hypothetical example. The issue is whether the estimate
describes a smooth attempt or a realistic range of attempts.

\section{Discussion}

What interests me is how easily an explanation for yesterday's delay can
become a reason to ignore it today. Each interruption may have been
unusual, yet a plan with no room for interruptions can still be too short.

I would be cautious about calling every missed deadline a planning
fallacy. A task can change after work begins. People can also choose an
ambitious deadline knowing that it is unlikely. A forecast and a target
answer different questions.

For a small project, I would trust a record of comparable work more than
the feeling that this attempt should be straightforward.

\section{Practical implications}

I would start with the time similar tasks actually took, then adjust
for differences I can explain. If the work is new, I would state the
uncertainty and revise the estimate after the first part is complete.

It also seems worth separating working time from elapsed time.
A draft that needs feedback may occupy only a few hours of attention
while remaining unfinished for several days.

\renewcommand{\refname}{References}
\begin{thebibliography}{9}
\small
\setlength{\itemsep}{0.35em}

\bibitem{buehler}
Buehler, R., Griffin, D., and Ross, M. (1994).
\href{https://doi.org/10.1037/0022-3514.67.3.366}
{Exploring the planning fallacy: Why people underestimate their task completion times.}
\textit{Journal of Personality and Social Psychology}, 67(3), 366--381.

\bibitem{kahneman}
Kahneman, D., and Lovallo, D. (1993).
\href{https://doi.org/10.1287/mnsc.39.1.17}
{Timid choices and bold forecasts: A cognitive perspective on risk taking.}
\textit{Management Science}, 39(1), 17--31.

\end{thebibliography}
\end{document}
